\documentclass[11pt]{article}

\usepackage[T1]{fontenc}
\usepackage[utf8]{inputenc}
\usepackage{lmodern}
\usepackage{amsmath,amssymb,amsthm}
\usepackage{xcolor}
\usepackage{hyperref}

\hypersetup{colorlinks=true, linkcolor=blue!50!black, urlcolor=blue!50!black}

\theoremstyle{definition}
\newtheorem{definition}{Definition}
\theoremstyle{plain}
\newtheorem{theorem}{Theorem}
\newtheorem{lemma}{Lemma}
\renewcommand{\thelemma}{\Alph{lemma}}

% The status badges of the website, as they appear next to a concept:
% a green "proven" chip and a yellow "open" chip.
\definecolor{provenbg}{HTML}{E8F5E8}
\definecolor{provenfg}{HTML}{2D6A2D}
\definecolor{provenborder}{HTML}{B8DAB8}
\definecolor{openbg}{HTML}{FFF6DC}
\definecolor{openfg}{HTML}{7A5A00}
\definecolor{openborder}{HTML}{E8D9A6}
\newcommand{\laxbadge}[4]{{\setlength{\fboxsep}{1.5pt}\setlength{\fboxrule}{0.4pt}%
  \fcolorbox{#3}{#2}{\textcolor{#1}{\scriptsize\sffamily\bfseries #4}}}}
\newcommand{\laxproven}{\laxbadge{provenfg}{provenbg}{provenborder}{THM\,\checkmark}}
\newcommand{\laxopen}{\laxbadge{openfg}{openbg}{openborder}{THM\,$\times$}}

\title{An Introduction to Lax}
\author{Jan Dreier}
\date{September 2026}

\begin{document}
\maketitle

\noindent
Lax is an archive that links natural-language mathematics with Lean
formalizations: an ``arXiv for Lean'', where submissions build on each other
towards a shared body of formalized mathematics.
% lax begin lax-242665
This document is itself a Lax submission, made from the features it describes.
% lax end
A submission consists of \emph{concepts}, \emph{proofs}, and optionally a
\emph{paper}, which we introduce in turn, using the prime numbers as a running
example.

\section{Concepts}

Lax represents mathematical definitions and claims as so-called
\emph{concepts}. A concept pairs a natural-language statement, as it would
appear in a paper, with a faithful encoding of that statement in Lean. This
section is annotated with two of them, shown on the right.

% lax begin Lax242665.Primes
\begin{definition}\label{def:prime}
  A natural number greater than $1$ is \emph{prime} if it is divisible only
  by $1$ and by itself.
\end{definition}
% lax end

% lax begin Lax242665.InfinitelyManyPrimes
\begin{theorem}[Euclid]\label{thm:euclid}
  For every natural number $n$ there is a prime $p > n$.
\end{theorem}
% lax end

Even if you are not familiar with Lean, try reading the Lean code to see
whether it encodes the intended meaning. One thing stands out:
Theorem~\ref{thm:euclid} is stated not as a \texttt{theorem} but as an
\texttt{axiom}, without a proof. That is Lax's way of separating claims from
proofs. The correctness of a formal proof is guaranteed by Lean, so with Lax,
readers only review the one thing Lean \emph{cannot} check: whether the Lean
code faithfully represents the intended mathematics. Concepts are therefore
written in especially clean and simple Lean code that even non-experts can
read. Still, mismatches can be subtle, so it is best to get many pairs of eyes
involved. You can click on each concept to see its full context, including a
discussion thread where the faithfulness of the encoding can be debated.

Concepts can be reused across submissions, which lets established academic
practices transfer to Lax: submissions build upon each other, and common
standards and definitions emerge by an organic selection process. For example,
this submission introduces prime numbers as a concept and deliberately leaves
one claim about them open (Lemma~\ref{lem:bertrand} below); a follow-up
submission can build on it, reuse Definition~\ref{def:prime} unchanged, and
supply the missing proof. The definition of a prime number is of course
already in mathlib~\cite{mathlib2020}. For notions that are not, such as
% lax begin Lax48.TwinWidth
twin-width% lax end
,
% lax begin Lax12.NowhereDenseClasses
nowhere dense graph classes% lax end
, or
% lax begin Lax67.Ram
the word RAM% lax end
, Lax can help the community agree on common standards.

\section{Proofs}

Let us now turn towards proofs, such as this one.

% lax begin Lax242665Proofs.InfinitelyManyPrimes.exists_prime_gt
\begin{proof}[Proof of Theorem~\ref{thm:euclid}]
  Let $p$ be the smallest prime factor of $n! + 1$. If we had $p \leq n$,
  then $p$ would divide $n!$, and hence also $(n! + 1) - n! = 1$, which is
  impossible. So $p > n$.
\end{proof}
% lax end

It is annotated on the right with the corresponding Lax proof. A
natural-language proof serves two purposes: it convinces the reader that a
claim holds, and it explains why. Lean proofs fully guarantee the former, but
are not particularly suited for the latter. They can run to enormous length
and are rarely enlightening even when short. Lax therefore treats formal proofs
as opaque blobs.

A green badge \laxproven{} next to a concept indicates that it has a complete
formal proof; a yellow one \laxopen{} indicates that it does not. However, Lax
does not ensure that the formal proof follows the same argument as the prose,
or even that the prose is correct. To gain more control over the argument,
authors may introduce concepts for the important lemmas along the way, thereby
certifying each step of the informal proof rather than only its conclusion.

Lax tracks which lemmas were used to prove which other lemmas.
Formally, each proof assumes a set of concepts and derives a target concept.
This defines a \emph{proof network} whose nodes are concepts and whose
hyperedges are proofs between them. For example, here
Lemmas~\ref{lem:odd}~and~\ref{lem:bertrand} are used to derive
Theorem~\ref{thm:oddprime}.

% lax begin Lax242665.OddPrimes
\begin{lemma}\label{lem:odd}
  Every prime other than $2$ is odd.
\end{lemma}
% lax end

% lax begin Lax242665Proofs.OddPrimes.odd_of_prime
\begin{proof}
  Let $p \neq 2$ be prime. If $p$ were even, then $2$ would divide $p$, so by
  Definition~\ref{def:prime} we would have $2 = 1$ or $2 = p$; both are
  impossible.
\end{proof}
% lax end

% lax begin Lax242665.BertrandPostulate
\begin{lemma}[Bertrand's postulate]\label{lem:bertrand}
  For every natural number $n \geq 1$ there is a prime $p$ with
  $n < p \leq 2n$.
\end{lemma}
% lax end

% lax begin Lax242665.OddPrimeBetween
\begin{theorem}\label{thm:oddprime}
  For every natural number $n \geq 2$ there is an odd prime $p$ with
  $n < p \leq 2n$.
\end{theorem}
% lax end

% lax begin Lax242665Proofs.OddPrimeBetween.exists_odd_prime_between
\begin{proof}
  By Lemma~\ref{lem:bertrand} there is a prime $p$ with $n < p \leq 2n$.
  Since $n \geq 2$, we have $p > 2$, so $p$ is odd by Lemma~\ref{lem:odd}.
\end{proof}
% lax end

Note that Lemma~\ref{lem:bertrand} is not proven in this submission. A
follow-up submission may provide a proof, thereby completing the proof of
Theorem~\ref{thm:oddprime}, after which both are marked as proven
\laxproven{} on the website. Clicking on a concept shows its proof network
and its open proof obligations. Providing a concept and proving it can thus be
separate contributions by separate people.

\section{Papers}

Besides concepts and proofs, a submission may carry a \LaTeX{} document, like
this one. Comments of the form \texttt{\% lax begin <id>} and
\texttt{\% lax end} annotate the text with Lax objects, where
\texttt{<id>} names a concept, a
proof, or a submission. A normal \LaTeX{} build ignores these comments, so
you may upload the same document to arXiv and to Lax. Using Rados\l{}aw
Pi\'orkowski's ReflowTeX library~\cite{reflowtex}, the \TeX{} source is
rendered for the web as well, reflowed to the reader's screen width beside
the as-printed PDF.

\section{Submissions}

Lean's mathematics traditionally lives in one curated library,
mathlib~\cite{mathlib2020}. Lax is a bet on an alternative model: that formal
mathematics can grow the way informal mathematics always has, as a literature
of independent papers that build on one another without a central authority.

For this to work, submissions must be stable enough to cite. As on arXiv, a
registered submission is immutable. This submission, with id
\texttt{lax-242665}, can be cited via
\begin{verbatim}
@misc{lax-242665,
  author       = {Dreier, Jan},
  title        = {An Introduction to Lax},
  year         = {2026},
  howpublished = {Lax, \url{https://laxarchive.org/lax-242665/}}
}
\end{verbatim}
This also makes it easy to attach a Lax submission to a conference or journal
paper under review; the author list may be left empty if the review process
requires anonymity. If changes are needed, you upload a revision under a new
id that \emph{supersedes} the old one, and the website points readers from the
old version to the new one.

Immutability has a cost on Lax that it does not have on arXiv. A Lean
submission can only build on another if both use the same version of mathlib,
and mathlib evolves over time. Lax therefore declares a \emph{baseline
version} of mathlib (currently \texttt{v4.30.0}). Submissions may use any
version and are verified and archived either way, but only submissions on the
same baseline version can build on each other. The baseline advances rarely.
When it does, an older submission is brought forward by a revision, typically
by whoever wants to build on it.

\section{Writing a submission}

The easiest way to get started is the \texttt{lax} command-line tool. Install
it with
\begin{verbatim}
npm install -g lax-archive && lax doctor
\end{verbatim}
and, inside a public git repository, run \texttt{lax init}. This creates
the layout the archive expects:
\begin{verbatim}
manifest.yaml     title, authors, environment, and what it builds on
abstract.md       the summary shown on the submission's page
main.tex          optional, the LaTeX document
LICENSE           Apache 2.0, the license of Lean and mathlib
concepts/         a Lean package holding the concepts, one module each
    lakefile.toml
    ...
proofs/           a second Lean package holding the proofs
    lakefile.toml
    ...
\end{verbatim}
Next, a prompt like ``run \texttt{lax print instructions} and formalize the
following result'' to your favorite coding agent will populate these folders
with concepts and proofs. Your agent will take over and guide you through
the whole process of creating a submission. This may take several sessions,
depending on the scope of the project. Your agent uses \texttt{lax build} to
make sure the project is accepted by the archive, and you can run
\texttt{lax serve} in the project folder to preview a rendering of the work
in progress. Once you and your agent are happy, run \texttt{lax submit},
which sends the name of the git repository (together with folder and commit
hash) to the archive, where it is validated and published. Initially, your
submission is in draft state and can be freely overwritten by further
submits. When you are ready, run \texttt{lax register} to make the submission
permanent and citable, so that others can build upon it.

\begin{thebibliography}{9}
\bibitem{mathlib2020}
  The mathlib Community.
  The Lean Mathematical Library.
  In \emph{Proceedings of the 9th ACM SIGPLAN International Conference on
  Certified Programs and Proofs (CPP 2020)}, pages 367--381, 2020.
\bibitem{reflowtex}
  Rados\l{}aw Pi\'orkowski.
  ReflowTeX: reflowable web rendering of \LaTeX{} sources.
  \url{https://github.com/radek-p/reflowtex}, 2026.
\end{thebibliography}

\end{document}
